\documentclass[10pt]{article}
\usepackage[T1]{fontenc}
\usepackage[utf8]{inputenc}
\usepackage[french]{babel}
\usepackage{amsmath,amssymb,mathrsfs,mathpazo}
\usepackage{geometry}
\geometry{a4paper,margin=2cm}
\usepackage{graphicx,array,tabularx,tabularray,colortbl,multirow,multicol,booktabs,enumitem,pifont,ulem,fancybox,diagbox,cancel}
\usepackage{pstricks,pstricks-add,pst-plot,pst-tree,pst-node,pst-text,pst-func,pst-math,pst-bspline,pst-3dplot,pst-eucl,pst-solides3d}
\let\pscalculate\undefined
\usepackage{pst-bezier}
\usepackage{tikz,pgfplots,tkz-tab}
\usetikzlibrary{arrows,arrows.meta,patterns,calc,positioning,trees,decorations.markings,shapes.geometric,shapes.misc,shapes.arrows,intersections,decorations.pathmorphing,decorations.pathreplacing,backgrounds,shadows,babel}
\pgfplotsset{compat=1.18}
\usepgfplotslibrary{fillbetween}
\usepackage[np]{numprint}
\usepackage{listings,makecell,pas-tableur,twoopt}
\let\barre\undefined
\usepackage[pstricks]{tablvar}
\definecolor{darkgreen}{RGB}{0,100,0}
\providecommand{\red}{\color{red}}
\providecommand{\blue}{\color{blue}}
\providecommand{\green}{\color{green}}
\providecommand{\black}{\color{black}}
\providecommand{\white}{\color{white}}
\providecommand{\gray}{\color{gray}}
\providecommand{\cyan}{\color{cyan}}
\providecommand{\magenta}{\color{magenta}}
\providecommand{\yellow}{\color{yellow}}
\providecommand{\violet}{\color{violet}}
\providecommand{\orange}{\color{orange}}
\providecommand{\euro}{\texteuro}
\providecommand{\EUR}{\texteuro}
\providecommand{\ieme}{\textsuperscript{e}}
\providecommand{\degres}{\ensuremath{^\circ}}
\providecommand{\degree}{\ensuremath{^\circ}}
\providecommand{\vv}[1]{\overrightarrow{#1}}
\providecommand{\Coord}[1]{#1}
\providecommand{\decosix}{}
\providecommand{\starredbullet}{\ensuremath{\bullet}}
\pagestyle{empty}
\setlength{\parindent}{0pt}
\usepackage{adjustbox}\DeclareUnicodeCharacter{2236}{:}\DeclareUnicodeCharacter{2190}{\ensuremath{\leftarrow}}
\providecommand{\R}{\mathbb{R}}
\providecommand{\N}{\mathbb{N}}
\providecommand{\Z}{\mathbb{Z}}
\providecommand{\Q}{\mathbb{Q}}
\providecommand{\D}{\mathbb{D}}
\providecommand{\C}{\mathbb{C}}
\providecommand{\e}{{\mathrm{e}}}
\providecommand{\vect}[1]{\overrightarrow{#1}}
\providecommand{\vectt}[1]{\overrightarrow{\mathrm{#1}}}
\providecommand{\vvt}[1]{\overrightarrow{\mathrm{#1}}}
\providecommand{\np}[2][]{{#2~#1}}
\providecommand{\vv}[1]{{\overrightarrow{#1}}}
\renewcommand{\d}{{\mathrm{d}}}
\providecommand{\degres}{{^\circ}}
\providecommand{\degre}{{^\circ}}
\providecommand{\blue}{\color{blue}}
\providecommand{\red}{\color{red}}
\providecommand{\enskip}{\quad}
\providecommand{\llbracket}{\lbrack\!\lbrack}
\providecommand{\rrbracket}{\rbrack\!\rbrack}
\providecommand{\euro}{{\text{€}}}
\providecommand{\no}{{\text{n°}}}
\providecommand{\interval}[2]{{[#1;#2]}}
\providecommand{\intervFO}[2]{{[#1;#2[}}
\providecommand{\dotfill}{\ldots}

\providecommand{\ds}{\displaystyle}
\providecommand{\textding}[1]{#1}
\providecommand{\Oij}{\(\mathrm{O};\vec{\imath},\vec{\jmath}\)}
\providecommand{\Oijk}{\(\mathrm{O};\vec{\imath},\vec{\jmath},\vec{k}\)}
\providecommand{\Ouv}{\(\mathrm{O};\vec{u},\vec{v}\)}
\providecommand{\intervFF}[2]{\ensuremath{[#1;#2]}}
\renewcommand{\euro}{\ifmmode\text{\texteuro}\else\texteuro\fi}
\DeclareUnicodeCharacter{00B0}{\ifmmode^{\circ}\else\textdegree\fi}
\DeclareUnicodeCharacter{20AC}{\ifmmode\text{\texteuro}\else\texteuro\fi}
% Énoncés : Éducation nationale. Transcription et corrigés : APMEP, Denis Vergès et contributeurs. Découpage et classement : CoopMaths, Sébastien Lozano, Éric Elter et contributeurs. Adaptation de la présentation : Maths Explorer.
% CC BY-SA : https://coopmaths.fr/www/about/
% https://coopmaths.fr/alea/static/bac/2021/tex/bac_2021_06_asie_9_cor.tex
\begin{document}
\noindent\textit{APMEP / CoopMaths. CC BY-SA. Présentation adaptée par Maths Explorer.}\par\medskip
\medskip

\begin{tabular}{|l|}\hline
\textbf{Principaux domaines abordés}\\
-- Suites\\
-- Équations différentielles\\ \hline
\end{tabular}







\bigskip

\textbf{Partie I : modèle discret}

\medskip









\medskip

\begin{enumerate}
\item 
On a pour $n = 0$, \, $u_1 = u_0 + 0,05(20 - u_0) = 1 + 0,05 \times 19 = 1 + 0,95 = 1,95$.
\item 
	\begin{enumerate}
		\item 
Pour tout naturel $n$, \, $u_{n+1} = u_n + 0,05\left(20 - u_n\right) = u_n  + 1 - 0,05u_n = u_n(1 - 0,05) + 1 = 0,95u_n + 1$.
		\item 
		

Pour tout naturel $n$, \, $v_{n+1} = 20 - u_{n+1} = 20 - \left(0,95u_n + 1\right) = 20 - 0,95u_n - 1 = $

$19 - 0,95u_n = 0,95 \times  20 - 0,95u_n = 0,95\left(20 - u_n \right) = 0,95v_n$.

Conclusion : pour tout naturel $n$,\, $v_{n+1} = 0,95v_n$ : cette égalité montre que la suite $\left(v_n\right)$ est une suite géométrique de terme initial $v_0 =  20 - u_0 =  20 - 1 = 19$ et de raison 0,95.
		\item 
On sait que pour tout $n \in n\N$, \, $v_n = v_0 \times q^n$, $q$ étant la raison, soit $v_n = 19  \times 0,95^n$.

Or $v_n = 20 - u_n \iff u_n = 20 - v_n = 20 - 19 \times 0,95^n$.
	\end{enumerate}
\item 
On vient de démontrer que pour tout naturel $n$, \, $u_n = 20 - 19\times 0,95^n$.

Comme $0 < 0,95 < 1$, on sait que $\displaystyle\lim_{n \to + \infty} 0,95^n = 0$, d'où par somme de limites :

\[\displaystyle\lim_{n \to + \infty} u_n = 20.\]
\end{enumerate}

\bigskip

\textbf{Partie II : modèle continu}

\medskip





\[(E) \qquad y' = 0,05(20 - y)\]









\begin{enumerate}
\item 
$L$ est la somme  de fonctions dérivables sur $[0~;~+\infty[$ et sur cet intervalle :

$L'(t) = - 0,05 \times \left(- 19\e^{-0,05t}\right) = 0,95\e^{-0,05t}$.

Donc $L$ est solution de $(E)$ si :

$y' = 0,05(20 - y) \iff 0,95\e^{-0,05t} = 0,05\left(20 - (\left(20 - 19\e^{-0,05t}\right)\right) \iff  0,95\e^{-0,05t}  = 0,05\left (19\e^{-0,05t}\right) \iff  0,95\e^{-0,05t}  =  0,95\e^{-0,05t}$ qui est vraie.

De plus $L(0) = 20 - 19\e^{-0,05\times 0} = 20 - 19\times 1 = 1$.
\item 
	\begin{enumerate}
		\item 
		
\starredbullet~$L'(0) = 0,95\e^{-0,05\times 0} =  0,95 \times 1 =  0,95$.

\starredbullet~$L'(5) = 0,95\e^{-0,05\times 5} =  0,95 \times \e^{-0,25}  \approx  0,74$.

Donc $L'(0) > L'(5)$.
		\item 
		

On sait que $\displaystyle\lim_{t \to + \infty} \e^{-0,05t} = 0$, donc $\displaystyle\lim_{t \to + \infty} L'(t) = 0$.

Ce résultat est bien en cohérence avec la description du modèle de croissance du bambou : celui-ci a une taille croissante ($L'(t) > 0$) de 1~m (taille initiale)  à 20~m (taille finale), la dérivée  donc la vitesse de croissance se rapprochant de zéro.
	\end{enumerate}
\end{enumerate}

\bigskip

\end{document}
